% ===================== CHAPITRE 3 — SPRINT 1 =========================
\chapter{Sprint 1 : Authentification et gestion des rôles}

\section*{Introduction}
Ce chapitre couvre le premier sprint de réalisation, consacré à la sécurité
d'accès. Après un rappel des objectifs, il présente le diagramme de cas
d'utilisation et le \emph{sprint backlog}, l'analyse des besoins, la conception
au moyen d'un diagramme de séquence, puis les interfaces réalisées.

\section{Objectifs du Sprint}
L'objectif de ce sprint est de mettre en place l'\textbf{authentification} des
utilisateurs et la \textbf{gestion des rôles et des droits d'accès}
(\texttt{ADMIN}, \texttt{DIRECTEUR\_COMMERCIAL}, \texttt{USER}), socle de sécurité
de toute l'application.

\section{Spécification des besoins}

\subsection{Diagramme de cas d'utilisation du Sprint 1}
La figure~\ref{fig:uc_s1} présente les cas d'utilisation du sprint.

\begin{figure}[htbp]
  \centering
  \includegraphics[width=0.82\textwidth]{uc_sprint1.png}
  \caption{Diagramme de cas d'utilisation du Sprint 1}
  \label{fig:uc_s1}
\end{figure}

\subsection{Sprint Backlog du sprint 1}
Le tableau~\ref{tab:bl_s1} détaille les user stories du sprint et les tâches
techniques associées.

\begin{table}[htbp]
\centering
\caption{Sprint Backlog du sprint 1}
\label{tab:bl_s1}
\renewcommand{\arraystretch}{1.3}
\begin{tabularx}{\textwidth}{>{\bfseries}p{1cm}p{5.2cm}X}
\toprule
\textmd{\textbf{US}} & \textbf{User story} & \textbf{Tâches techniques}\\
\midrule
US1 & S'authentifier avec identifiant et mot de passe. &
Créer le \texttt{LoginServlet} (\texttt{POST /login}) ; vérifier les identifiants
via le \texttt{SecurityService} sur la table \texttt{SECURITY\_USERS} ; ouvrir une
\texttt{HttpSession} ; créer le \texttt{LogoutServlet}.\\
\midrule
US2 & Gérer les rôles et les droits d'accès. &
Associer un rôle (\texttt{ADMIN}, \texttt{DIRECTEUR\_COMMERCIAL}, \texttt{USER}) à
chaque utilisateur ; restreindre l'accès aux fonctionnalités selon le rôle ;
gérer la réinitialisation du mot de passe (\texttt{ForgotPasswordServlet}).\\
\bottomrule
\end{tabularx}
\end{table}

\section{Analyse des besoins}

\subsection{Description textuelle des cas d'utilisation}
Les tableaux~\ref{tab:cu_auth} et~\ref{tab:cu_roles} décrivent les cas
d'utilisation principaux du sprint.

\cudesc{tab:cu_auth}{S'authentifier}
{Permettre à l'utilisateur de se connecter à l'application avec son identifiant et
son mot de passe.}
{Administrateur / Directeur commercial / Utilisateur}
{L'utilisateur doit disposer d'un compte actif dans la table \texttt{SECURITY\_USERS}.}
{Une session est ouverte et l'utilisateur est redirigé vers le tableau de bord
correspondant à son rôle.}
{(1) L'utilisateur saisit son identifiant et son mot de passe ; (2) le
\texttt{SecurityService} vérifie les informations sur la table de sécurité ; (3) en
cas de succès, une session est ouverte et l'utilisateur est redirigé ; (4) sinon,
un message d'erreur est affiché.}

\cudesc{tab:cu_roles}{Gérer les rôles et les droits d'accès}
{Permettre à l'administrateur d'attribuer un rôle à chaque utilisateur afin de
contrôler l'accès aux fonctionnalités.}
{Administrateur}
{L'administrateur doit être authentifié avec le rôle \texttt{ADMIN}.}
{Le rôle est enregistré et les droits d'accès de l'utilisateur sont mis à jour.}
{(1) L'administrateur consulte la liste des utilisateurs ; (2) il sélectionne un
utilisateur ; (3) il lui attribue un rôle (\texttt{ADMIN},
\texttt{DIRECTEUR\_COMMERCIAL} ou \texttt{USER}) ; (4) le système applique le
contrôle d'accès correspondant.}

\section{Conception}

\subsection{Diagramme de séquence : « S'authentifier »}
La figure~\ref{fig:seq_auth} décrit le déroulement dynamique de l'authentification.

\begin{figure}[htbp]
  \centering
  \includegraphics[width=0.92\textwidth]{seq_auth.png}
  \caption{Diagramme de séquence -- Authentification}
  \label{fig:seq_auth}
\end{figure}

Le \texttt{LoginServlet} reçoit la requête \texttt{POST /login} et délègue la
vérification au \texttt{SecurityService}, qui interroge la table
\texttt{SECURITY\_USERS} via l'\texttt{EntityManager}. Le fragment combiné
\texttt{alt} distingue les deux issues : en cas de succès, l'\texttt{Utilisateur}
métier est associé (ou créé), une session est ouverte et l'utilisateur est
redirigé vers le tableau de bord ; en cas d'échec, un message d'erreur est
renvoyé.

\section{Réalisation : interfaces utilisateur}
La figure~\ref{fig:ui_login} présente la page d'authentification réalisée :
l'utilisateur y saisit son identifiant et son mot de passe pour accéder à
l'application.

\begin{figure}[htbp]
  \centering
  \includegraphics[width=\textwidth]{screens/login.png}
  \caption{Interface d'authentification}
  \label{fig:ui_login}
\end{figure}

\section*{Conclusion}
Ce sprint a mis en place le socle de sécurité de l'application :
l'authentification et la gestion des rôles. Le sprint suivant aborde la gestion
des agences et des employés.
